\section{Results and Discussion}

\subsection{\texorpdfstring{${[\ce{NO2}]_{\mathrm{equiv}}}$ in the Troposphere Determined from Kinetic Data}{[NO2]equiv in the Troposphere Determined from Kinetic Data}}
Figure \ref{DudleySurfPlot} illustrates that the ${[\ce{NO2}]_{\mathrm{equiv}}}$ determined from kinetic data decreases with temperature and altitude and that the most likely place where the equivalent point may be reached is in fact around the cruise altitude of long-haul aircraft. Short-haul aircraft may not be able to attain such altitudes if the flight is very short in duration. The simple model presented in figure \ref{DudleySurfPlot} suggests that as the flight reaches cruise altitude at about 10~km, the amount of \ce{NO2} required for parity in OH loss is around 2~ppb. How feasible is it to generate such levels of \ce{NO2}? In a narrow flight corridor with frequent flights through it, it may well be.


\subsection{\texorpdfstring{Global Mapping of ${[\ce{NO2}]_{\mathrm{equiv}}}$ at the Cruise Level}{Global Mapping of [NO2]equiv at the Cruise Level}}
Figure~\ref{NOxt_stochem} shows global ${[\ce{NO2}]_{\mathrm{equiv}}}$ at model output level 8 which is determined by using [\ce{CH4}], [CO], [M], and temperature fields simulated by STOCHEM-CRI. In this study, data at model output level 8 (9.2--11.8 km) were used, to best replicate the conditions experienced by civil aircraft in cruise. It is known that CO mixing ratios are very variable in space and time, for example, in the biomass burning season in the tropics, CO in the upper troposphere can be elevated considerably, which could be responsible for the spatial and temporal variations of global ${[\ce{NO2}]_{\mathrm{equiv}}}$ level. 

This figure suggests that sub-ppb levels of ${[\ce{NO2}]_{\mathrm{equiv}}}$ may exist in mid- and high-latitudes, but that several ppb (up to 2~ppb) are required in the tropics, where temperature is higher, and CO levels elevated through deep convection of CO from the surface sources (e.g., vegetation, biomass burning, anthropogenic emissions, VOC oxidation). The biomass burning season for Northern Hemisphere Africa and South East Asia is from November to March with peak burning in December and January; however, the biomass burning seasons for South America and Southern Africa show a peak in September to October with little burning between November and February, which is responsible for the monthly variations of ${[\ce{NO2}]_{\mathrm{equiv}}}$ peaks at these areas (see Figure \ref{NOxt_stochem}). There is little biomass burning over Central Europe and North America, but the strong anthropogenic source along with VOC oxidation leads to a higher level of CO in these areas, resulting in \textasciitilde1~ppb ${[\ce{NO2}]_{\mathrm{equiv}}}$. The vegetation emission of CO can also increase ${[\ce{NO2}]_{\mathrm{equiv}}}$ at Amazon and South Africa forest regions during Austral summer months (December-January-February).

To further confirm the \ce{NO_x} threshold, termed ${[\ce{NO2}]_{\mathrm{equiv}}}$, trends seen in the global model analysis in Figure \ref{NOxt_stochem}, the results using IAGOS data are plotted in Figure \ref{NOxt_iagos}. The ${[\ce{NO2}]_{\mathrm{equiv}}}$ values obtained from both STOCHEM and IAGOS show a good agreement (see Figures \ref{NOxt_stochem} and \ref{NOxt_iagos}) with the differences between all months averaged ${[\ce{NO2}]_{\mathrm{equiv}}}$ at level 8 referred to by mean bias (STOCHEM-IAGOS) and the percentage change ((STOCHEM-IAGOS)*100/IAGOS) are found to be 0.147~ppb and 0.27\%, respectively. The calculated global ${[\ce{NO2}]_{\mathrm{equiv}}}$ from both model and measurement data (Figures \ref{DudleySurfPlot} and \ref{NOxt_stochem}) suggested that the northern midlatitudes of the atmosphere are the most favourable region existing with the smallest \ce{NO_x} thresholds (0.5~ppb) needed before reduction in ozone formation is likely to occur at aircraft cruise altitudes.

\begin{figure}[H]
		\includegraphics[width=\textwidth]{4/NOxt_stochem.jpg}
		%\vspace{-.1cm}
		\caption{Global ${[\ce{NO2}]_{\mathrm{equiv}}}$ at model output level 8 (9.2--11.8 km) estimated using [\ce{CH4}], [CO], [M], and temperature fields simulated by STOCHEM Common Representative Intermediate (CRI).}
		\label{NOxt_stochem}
\end{figure}

\begin{figure}[H]
		\includegraphics[width=\textwidth]{4/NOxt_iagos.jpg}
		%\vspace{-.1cm}
		\caption{Global ${[\ce{NO2}]_{\mathrm{equiv}}}$ at model output level 8 (9.2--11.8 km) estimated using [CO], [M], and temperature fields from In-service Aircraft for Global Observing System (IAGOS) dataset and [\ce{CH4}] from STOCHEM Common Representative Intermediate (CRI) model simulation.}
		\label{NOxt_iagos}
\end{figure}

Assessing how well \ce{NO_x} is sequestered by OH radicals to form \ce{HNO3}, terminating ozone production, will depend on $k_6$ (function of pressure and temperature) and this itself will change with altitude, latitude, and time of year. As altitude increases for a given location, temperature decreases and so $k_6$ will increase. However, at the same time pressure is decreasing and so $k_6$ will decrease, determining which process wins and whether $k_6$ ultimately decreases or increases with altitude requires high fidelity kinetic information. An increasing $k_6$ with altitude may mean that this termination reaction is more effective at high altitudes for a given [\ce{NO2}]. These reactions are hard to study at low temperature and as a function of pressure and so uncertainties increase. The rate coefficients ($k_1$, $k_6$,and $k_7$) were perturbed upwards and downwards by considering the uncertainties of the rate coefficients and used to calculate global annual average ${[\ce{NO2}]_{\mathrm{equiv}}}$ (see Supporting Information figure S1 \url{https://onlinelibrary.wiley.com/doi/full/10.1002/kin.21644}). The uncertainty analysis shows that ${[\ce{NO2}]_{\mathrm{equiv}}}$ can be decreased and increased by up to 2.5\% and 5\% in the northern midlatitudes due to the using lower limit and higher limit rate coefficients, respectively.

The calculation of ${[\ce{NO2}]_{\mathrm{equiv}}}$ did not consider VOCs contributions in the study. The methane mixing ratios have been rising in the troposphere since ca. 2010 but most VOCs are decreasing. For ethane and propane \cite{Li2022}, two of the longest lived VOCs, and species that will have elevated concentrations in the upper troposphere, current measured concentrations are indeed decreasing \cite{Li2022}, and suggest that these VOCs would contribute less than 5\% of the OH loss compared with \ce{CH4}. Given that loss of OH with \ce{CH4} is typically 30\% of that for CO, VOCs other than CO and \ce{CH4} can be ignored in any initial analysis such as carried out here. Even for acetone, another long-lived species with elevated levels in the upper troposphere higher levels are seen in the summer months (800~ppt) but drop to around 200~ppt in winter\cite{DuFour2016}. So, for short periods of time in a year, levels approach that required to be around 5\% of that for \ce{CH4} but for the Northern Hemisphere, where it is expected the formation flying to be most effective, levels are much lower.

\subsection{Aircraft Plume and its Chemistry Study using a 1D Dispersion Model}
In a first box model experiment, an aircraft emits a pulse of \ce{NO_x} into the centre box of the model grid at 12:00~z. The pulse size was set based on a global aircraft cruise \ce{NO_x} emission of 3.9~Tg \ce{NO_x} per year and a distance travelled \cite{Wasiuk2016} of 36.4~billion~km that is 107~g \ce{NO_x}/km. This \ce{NO_x} was instantaneously mixed into the centre box volume and the time development of the model species was then followed. After several seconds, the NO-\ce{NO2}-\ce{O3} chemistry was rapidly established, a strong depletion in ozone formed and then \ce{NO_x} began to disperse horizontally into the neighbouring boxes over the period of hours and days. After 1~h the horizontal distribution of \ce{NO_x} showed an aircraft plume with a width of 3~km, as shown in Figure \ref{Disp1}, which expanded to 6.8~km after 6~h and 13.1~km after 24~h. Plume width expanded with the square root of lapsed time, see Figure \ref{Disp2}, exactly as predicted by the Einstein-Smoluchowski equation.

\begin{figure}[h!]
		\centering
		\includegraphics[width=0.7\textwidth]{4/Disp1.jpg}
		%\vspace{-.1cm}
		\caption{\ce{NO_x} mixing ratios as a function of distance from the flight path at increasing times after the passage of an aircraft, each offset by 3~ppb.}
		\label{Disp1}
\end{figure}

\begin{figure}[h!]
		\centering
		\includegraphics[width=0.7\textwidth]{4/Disp2.jpg}
		%\vspace{-.1cm}
		\caption{Aircraft plume width as a function of travel time.}
		\label{Disp2}
\end{figure}

Initially, all of the ozone would have been completely wiped out in the centre of the flight-path by its rapid reaction with NO:

\reaction[NO+O3]{NO + O_3 -> NO_2 + O_2}

However, during daytime, the product \ce{NO2} is rapidly photolysed to form ozone. Overall, the reaction sequence \eqref{HO2+NO_4}--\eqref{O+O2+M_4} rapidly returns the mixture of NO-\ce{NO2}-\ce{O3} to a steady state in a few minutes in which the ozone levels are largely restored to their values before the passage of the aircraft.

The distribution of ozone perpendicular to the flight path is shown in Figure \ref{Disp3} for various times after the passage of the aircraft. One hour after the aircraft passage, ozone levels show the presence of a slight dip close to the centre of the aircraft plume demonstrating that the stationary state in the NO-\ce{NO2}-\ce{O3} system has been re-established. After 3~h, the central dip is still evident, but photochemical ozone production is starting to appear in the wings of the aircraft plume. As time passes, the production in the wings becomes more evident and the central dip is filled in. After 24~h, photochemical ozone production is well on its way and an ozone plume becomes apparent with a width of about 16~km and a height of just over 0.4~ppb at its centre. At this stage, the ozone plume is about 3~km wider than the \ce{NO_x} plume as horizontal dispersion moves the photochemically produced ozone into the background atmosphere. After 4~days, see Figure \ref{Disp3}, the plume is about 30~km wide and has a height of about 2.7~ppb.

\begin{figure}[h!]
		\centering
		\includegraphics[width=0.8\textwidth]{4/Disp3.jpg}
		%\vspace{-.1cm}
		\caption{\ce{O3} mixing ratios as a function of distance from the flight path at increasing times after the passage of an aircraft, each offset by 3~ppb.}
		\label{Disp3}
\end{figure}

In the next application of the aircraft plume dispersion model, the case of five aircraft flying at the same time but separated in space are considered. To begin with the five aircraft flying together without any spatial separation are reported and then further cases where they are separated by 100~m, 500~m, and so on up to 10~km are investigated. Figure \ref{Disp4} shows the time development over 5~days of the 99.9~km grid average \ce{O3} for a base case with one aircraft flying, five aircraft flying together and five aircraft with spatial separations increasing from 100~m to 10~km.

In each plot, the grid average \ce{O3} perpendicular to the flight path shows a steady increase over the 5~days. This increase has small depressions during nighttime when the NO-\ce{NO2}-\ce{O3} system adjusts to the lack of sunlight. The steady increase is driven by photochemical ozone production because the chosen initial conditions for these experiments have NO levels above the local \ce{NO_x} compensation point. Initially the plots for the five aircraft lie below the one aircraft case because of the more intense \ce{O3} depletion found with more aircraft. However, the situation reverses after 1~day and \ce{O3} levels in the five aircraft cases then lie above the one aircraft case. After further time has elapsed, the five aircraft cases begin to diverge with the cases with wider separations diverging the greatest.

\begin{figure}[h!]
		\centering
		\includegraphics[width=0.7\textwidth]{4/Disp4.jpg}
		%\vspace{-.1cm}
		\caption{Time development of the average \ce{O3} mixing ratios across the 99.9~km grid for a base case with one aircraft, five aircraft flying together and flying with horizontal separations increasing from 100~m to 10~km.}
		\label{Disp4}
\end{figure}

After 5~days, the case with the 10~km separation has a grid averaged \ce{O3} mixing ratio which is 1.2~ppb higher than the no-separation case, see Figure \ref{Disp4}. This excess \ce{O3} drops to just under 0.8~ppb with a horizontal separation of 5~km and to 0.2~ppb with a separation of 2~km. On this basis, the impact of aircraft emissions on spatially-averaged \ce{O3} could be halved if the aircraft could maintain separations inside 4~km.

Clearly, this is a first illustrative example and our aircraft plume dispersion model is a highly simplified representation of the complex interplay between aircraft cruise and the local environment in an air traffic corridor. However, the results are encouraging enough for the continuation of model development, increasing spatial resolution, examining more complex scenarios, and assessing further proposals for the formation flying of aircraft.

